\documentclass[11pt,a4paper]{article}

\usepackage[utf8]{inputenc}
\usepackage[T1]{fontenc}
\usepackage[spanish]{babel}
\usepackage[a4paper,margin=2.3cm]{geometry}
\usepackage{array}
\usepackage{longtable}
\usepackage{enumitem}
\usepackage{xcolor}
\usepackage{listings}
\usepackage{hyperref}
\usepackage{parskip}

\hypersetup{
  colorlinks=true,
  linkcolor=black,
  urlcolor=blue!45!black,
  pdftitle={Arturito --- Informe de avance},
  pdfauthor={Equipo Arturito}
}

\urlstyle{tt}
\setlist{itemsep=2pt,topsep=4pt}

\newcommand{\ruta}[1]{\path{#1}}
\newcommand{\decision}[1]{\textbf{Decisi\'on:} #1}
\newcommand{\desvio}[1]{\textbf{Desv\'io respecto de MAIN:} #1}

\title{Arturito --- Informe de avance\\
       \large Documento vivo de ingenier\'ia (alcance, estado, plan)}
\author{Equipo Arturito \\
        Joaqu\'in Guzm\'an \and Tom\'as Gamarra \and Robaldi Santiago \and Goncalves Federico}
\date{3 de octubre de 2026}

\begin{document}
\maketitle

\begin{abstract}
Este archivo es el \textbf{documento de ingenier\'ia vivo} del rover Arturito.
Integra el objetivo acad\'emico de \ruta{ARTURITO MAIN.tex}, el estado real del
workspace y las decisiones de firmware acordadas para poder implementar sin
contradecirse. Cuando este documento y el c\'odigo difieran, hay que
\textbf{alinear el c\'odigo} (o actualizar este archivo en el mismo cambio).
El informe de c\'atedra \ruta{ARTURITO MAIN.tex} queda \textbf{congelado}: los
desv\'ios se registran aqu\'i, no se reescribe MAIN. El dise\~no de firmware
en capas se actualiza en un paso posterior; aqu\'i se anota en qu\'e no
coincide. \ruta{CONTEXTO\_PROYECTO\_ROVER.tex} queda \textbf{absorbido} por
este informe y deja de ser fuente de verdad.
\end{abstract}

\tableofcontents
\newpage

%==============================================================================
\section{C\'omo usar este documento}
%==============================================================================

\begin{itemize}
  \item \textbf{Antes} de una funci\'on nueva o un cambio de arquitectura:
        leer la secci\'on de decisiones congeladas y el plan.
  \item \textbf{Despu\'es} de implementarla: actualizar \emph{Estado},
        \emph{Registro de avances} y, si cambia un contrato, \emph{Decisiones}.
  \item No documentar experimentos personales que no definan al robot del grupo.
  \item Autoridad: \textbf{este archivo + el c\'odigo}. MAIN es entrega de
        c\'atedra (congelado). El \texttt{.tex} de dise\~no de firmware es
        meta a revisar, no contrato actual.
\end{itemize}

%==============================================================================
\section{Qu\'e es el proyecto (resumen de MAIN)}
%==============================================================================

Arturito es un \textbf{rover diferencial aut\'onomo mapeador} (Taller de
Proyecto~I, UNLP). Debe explorar un interior sin marcas en el piso, construir
una representaci\'on~2D y poder navegarla.

Problemas simult\'aneos: percepci\'on, localizaci\'on y planificaci\'on.
MAIN formula el horizonte como grilla de ocupaci\'on (y SLAM como
optimizaci\'on secundaria), con tracci\'on diferencial, encoders, IMU,
sensor de distancia, comunicaciones inal\'ambricas y GUI externa.

Este informe no reemplaza MAIN como entrega de c\'atedra. S\'i define
\textbf{qu\'e vamos a construir de verdad en el firmware} y en qu\'e orden.

%==============================================================================
\section{Decisiones congeladas (acuerdo de ingenier\'ia)}
%==============================================================================

Estas decisiones rigen el c\'odigo hasta que se cambien expl\'icitamente
en una actualizaci\'on de este archivo.

\subsection{Plataforma}

\decision{Desarrollar y cerrar capas en \textbf{ESP32} primero.
La EDU-CIAA se porta despu\'es reutilizando drivers y aplicaci\'on,
cambiando solo la HAL.}

MAIN no fija MCU. El dise\~no de firmware habla de FreeRTOS sobre EDU-CIAA.
ESP32 y EDU-CIAA ejecutan el mismo runtime FreeRTOS compartido. Cada placa
conserva su adaptador de arranque: Arduino invoca \texttt{setup()} con el
scheduler activo; CIAA entra por \texttt{main()}, crea el bootstrap y arranca
el scheduler.

\subsection{Scan de distancia (V0 / V1)}

\decision{ToF \textbf{fijo al chasis}. El robot gira sobre s\'i; el
\'angulo de cada medida es el yaw $\theta$ de la IMU, no un mecanismo
giratorio tipo LiDAR.}

\desvio{MAIN exige un mecanismo de rotaci\'on $360^\circ$ con posici\'on
angular propia del sensor. Eso no es el hardware ni el m\'etodo de esta etapa.}

\subsection{Encoders y motores (l\'inea del grupo)}

\decision{Encoders absolutos \textbf{AS5600} por rueda. Puente~H
\textbf{DRV8833}. El integrante a cargo de ruedas/encoders es Federico.
Santiago no define ni stub-ea esa API hasta que haya algo que medir.}

\desvio{MAIN habla de encoders por pulsos/giros en un intervalo. El
contrato de software ser\'a \'angulo absoluto AS5600, no cuenta de pulsos.}

Hasta que existan AS5600, la \textbf{pose oficial del grupo} para visualizaci\'on
V0 es
\[
q = [x, y, \theta]^\top = [0, 0, \theta_{\mathrm{IMU}}]^\top.
\]
No hay odometr\'ia de traslaci\'on oficial sin ruedas.

\subsection{Mapa: dos etapas}

\decision{Primero \textbf{V0}: nube / plot polar con el robot en el origen
(gira, no traduce). Despu\'es \textbf{V1}: grilla de ocupaci\'on~2D en marco
mundo, usando pose $(x,y,\theta)$ cuando existan AS5600. SLAM sigue siendo
objetivo secundario de MAIN, no el entregable de firmware actual.}

El grupo todav\'ia distingue mal ``plot de Processing'', ``grilla'' y ``SLAM''.
Este p\'arrafo es la definici\'on operativa para no mezclarlos.

\subsection{Capas de software (contrato actual)}

\decision{Tres capas. La capa~2 pide n\'umeros a funciones de capa~1 y
\textbf{no sabe} si debajo hay ESP32 o CIAA.}

Corte acordado (\textbf{HAL de perif\'ericos}, no HAL de buses):

\begin{itemize}
  \item \textbf{Capa 1 (HAL), una implementaci\'on por placa:} inicializar
        el perif\'erico y entregar lecturas ya en unidades de trabajo
        (p.\ ej.\ $\omega_z$ en rad/s, distancia en metros, m\'as adelante
        PWM/GPIO de motores). No calibra, no integra pose, no arma el mapa.
  \item \textbf{Capa 2 (drivers, en \ruta{shared}):} f\'isica y protocolo
        del chip a nivel de aplicaci\'on de sensor/actuador: offset y zona
        muerta de MPU, integraci\'on de $\theta$, validaci\'on de rango ToF;
        m\'as adelante DRV8833 y AS5600. Convierte a magnitudes usables por
        la app. No incluye \texttt{Arduino.h} ni SAPI.
  \item \textbf{Capa 3 (aplicaci\'on):} cinem\'atica / pose, grilla, navegaci\'on.
        Contiene el runtime FreeRTOS com\'un (adquisici\'on, drivers y
        publicaci\'on). La pose XY y la grilla siguen pendientes de encoders.
\end{itemize}

Esto \textbf{no} es el contrato de ejemplos \texttt{hal\_i2c\_read\_bytes()}
del dise\~no de firmware. Ese .tex se alinear\'a despu\'es.

\subsection{IMU en la l\'inea del grupo}

\decision{Validar y usar \textbf{giroscopio Z} $\rightarrow$ $\theta$.
MAIN pide aceleraci\'on en 3 ejes; no es requisito de esta etapa de grupo.}

\subsection{Visualizaci\'on y radio}

\decision{V0: Processing por \textbf{UART/USB}. Pronto (avance siguiente,
no este): telemetr\'ia de pose/mapa por Wi-Fi o Bluetooth del ESP32.}

\desvio{MAIN pone el enlace inal\'ambrico como objetivo primario de
comunicaciones. En ingenier\'ia, el primario actual es serie; el radio es
el pr\'oximo salto de comunicaciones.}

\subsection{FreeRTOS compartido}

\decision{Ambas placas ejecutan las mismas tres tareas desde
\ruta{shared/layer3\_app}: \textbf{adquirir} (HAL),
\textbf{actualizar drivers} ($\theta$ y distancia) y
\textbf{publicar} (telemetr\'ia).}

El scheduler ya est\'a activo cuando corre el bootstrap com\'un. ESP32 lo
inicia el framework Arduino; CIAA crea el bootstrap antes de llamar a
\texttt{vTaskStartScheduler()}. El tama\~no de stack se expresa en bytes en
la aplicaci\'on y se convierte a palabras para el FreeRTOS est\'andar de CIAA.

\subsection{Protocolo serie V0}

\decision{115200 baud. L\'ineas de arranque o log con prefijo \texttt{\#}
(Processing las ignora). Datos:}
\begin{lstlisting}
theta_rad,distancia_m
\end{lstlisting}
Distancia inv\'alida: \texttt{-1}. Una muestra por l\'inea, sin tercera columna.

Los sketches que esperan tres columnas o \texttt{distancia\_cm} se consideran
desactualizados: se adaptan al firmware, no al rev\'es.

\subsection{Organizaci\'on de archivos (estructura actual)}

HAL partida por perif\'erico para aislar fallas (p.\ ej.\ IMU sin ToF).
\texttt{BOARD\_*} solo desde el build, no \texttt{\#define BOARD\_EDUCIAA}
incondicional en \ruta{shared/config/config.h}.

\'Arbol actual del workspace:

\begin{lstlisting}
docs/                  % informes y documentos de arquitectura
shared/
      config/
  layer1_hal/
    hal_time.h
    hal_imu.h
    hal_tof.h
  layer2_drivers/
    driver_mpu6050.c/.h
    driver_distancia.c/.h
      layer3_app/
            app_runtime.c/.h      % tareas FreeRTOS comunes
esp32/src/
      main.cpp             % arranque Arduino; delega al runtime compartido
  layer1_hal/
    hal_time_esp32.cpp
    hal_imu_esp32.cpp
    hal_tof_esp32.cpp
            hal_telemetry_esp32.cpp
      tests/                   % diagnosticos: IMU sola / ToF solo
edu-ciaa/rover/src/layer1_hal/
      hal_time_educiaa.cpp
      hal_imu_educiaa.cpp
      hal_tof_educiaa.cpp
      hal_telemetry_educiaa.cpp
processing/
\end{lstlisting}

Init de IMU y ToF \textbf{separados}: si falta el VL53L0X,
\texttt{hal\_tof\_init()} informa el fallo y la IMU sigue publicando $\theta$.

%==============================================================================
\section{Estado actual del workspace (al 2026-10-03)}
%==============================================================================

Relevado a partir del c\'odigo y de lo que era
\ruta{CONTEXTO\_PROYECTO\_ROVER.tex}.

\begin{longtable}{>{\raggedright\arraybackslash}p{0.32\textwidth}
                  >{\raggedright\arraybackslash}p{0.60\textwidth}}
\textbf{Componente} & \textbf{Estado} \\
\hline
Workspace & Carpetas \ruta{esp32/}, \ruta{edu-ciaa/}, \ruta{shared/},
\ruta{processing/} y \ruta{docs/}; las cinco est\'an registradas en el
\texttt{.code-workspace}. \\
Capa 1 ESP32 & HAL separada en tiempo, IMU y ToF; la ausencia del ToF no
detiene la publicaci\'on de yaw. La implementaci\'on combinada heredada fue
retirada. \\
Capa 1 CIAA & HAL separada en tiempo, IMU y ToF; compilada, sin prueba en
placa. La falla del ToF no detiene la lectura de IMU. \\
Capa 2 & MPU: offset, zona muerta $1^\circ$/s, integraci\'on trapezoidal,
\texttt{get\_theta\_rad()} con convenci\'on de signo. Distancia: v\'alida
0.05--1.50~m, si no $-1$. \\
Capa 3 & \ruta{shared/layer3\_app/} contiene el runtime FreeRTOS com\'un;
pose y mapeo V1 quedan pendientes. Se retiraron \ruta{layer3\_robotics/} y
\ruta{layer4\_app/}. \\
ESP32 / CIAA & Tres tareas FreeRTOS compartidas (adquirir, drivers, publicar),
mutex y CSV V0 \texttt{theta\_rad,distancia\_m}; compiladas en ambas placas. \\
Motores / AS5600 & No hay drivers. \\
Mapa / nav / SLAM & No hay grilla ni pose XY. Processing
\texttt{flecha\_y\_mapa\_v1} fue validado para el plot polar V0. \\
\texttt{config.h} & No define placas; PlatformIO define
\texttt{BOARD\_ESP32}. \\
\end{longtable}

Limitaci\'on conocida (documentada, no es bug a ``arreglar'' con magia):
$\theta$ sale de integrar un gyro. Offset y zona muerta reducen deriva;
no la eliminan. No hay fusi\'on con aceler\'ometro en la l\'inea de grupo.

%==============================================================================
\section{Incongruencias abiertas (qu\'e no cierra)}
%==============================================================================

\subsection{Respecto de MAIN (congelado: se registran, no se editan all\'i)}

\begin{enumerate}
  \item Mecanismo LiDAR $360^\circ$ vs ToF fijo + yaw del chasis.
  \item Encoders por pulsos vs AS5600.
  \item IMU: accel 3 ejes como requisito de odometr\'ia vs solo gyro~Z.
  \item Comunicaciones inal\'ambricas primarias vs UART V0.
  \item Placeholders no funcionales (\texttt{[FreqMinMapeo]}, velocidades,
        tama\~no de mapa, etc.) nunca numerados.
  \item ``SLAM'' en la introducci\'on vs grilla V1 como meta de software
        y SLAM como secundario.
\end{enumerate}

\subsection{Respecto del dise\~no de firmware (.tex a actualizar despu\'es)}

\begin{enumerate}
  \item El dise\~no preliminar presenta EDU-CIAA como plataforma inicial;
        el workspace implementa ESP32 primero y un runtime FreeRTOS com\'un.
  \item HAL de buses (\texttt{hal\_i2c\_read\_bytes}) vs HAL de perif\'ericos
        acordada.
  \item Ejemplos de API de capa~3 y drivers de motor/encoder: meta, no c\'odigo.
\end{enumerate}

\subsection{Respecto del c\'odigo (ESP32 y compilaci\'on EDU-CIAA)}

El firmware V0 fue cargado y validado por Santiago con Processing. ESP32 y
EDU-CIAA comparten las tres tareas y drivers; ambos builds pasan. Resta probar
el comportamiento de la EDU-CIAA con sensores conectados.

\begin{enumerate}
  \item Hecho: HAL ESP32 separada (tiempo / IMU / ToF); init independiente.
  \item Hecho: entornos aislados \texttt{imu\_only} y \texttt{tof\_only}.
      \item Hecho: tareas FreeRTOS comunes (adquirir / drivers / publicar).
  \item Hecho: CSV V0 en metros, logs \texttt{\#} y parser Processing alineado.
  \item Hecho: HAL EDU-CIAA separada e integrada al build.
  \item Hecho: drivers compartidos usados por ambos builds.
  \item Pendiente: prueba de sensores en hardware EDU-CIAA.
  \item Hecho: workspace con tres capas compartidas; documentaci\'on en
        \ruta{docs/}; Processing visible en el workspace.
\end{enumerate}

%==============================================================================
\section{Plan de acci\'on}
%==============================================================================

\subsection{Paso 0 --- Ahora}

Contrato actualizado; ESP32 y EDU-CIAA crean el mismo runtime FreeRTOS y sus
builds pasan. El plot V0 fue validado en Processing. La EDU-CIAA requiere
prueba f\'isica de sensores antes de dar por cerrado ese port.

\subsection{Paso 1 --- Runtime FreeRTOS com\'un (implementado; build OK)}

Orden sugerido:

\begin{enumerate}
  \item Contratos \ruta{hal\_time.h}, \ruta{hal\_imu.h}, \ruta{hal\_tof.h}
        y splits \texttt{\_esp32.cpp}.
  \item Init IMU independiente del ToF.
  \item Drivers shared sin cambio de algoritmo (solo adaptaci\'on a los
        nuevos nombres / metros).
  \item \texttt{app\_runtime.c}: tres tareas compartidas; cada \texttt{main}
        queda como adaptador de arranque y plataforma.
  \item Test IMU-only (y ToF-only si hay hardware).
  \item Un sketch Processing que consuma exactamente el CSV V0.
\end{enumerate}

\subsection{Paso 2 --- Comunicaciones (pr\'oximo avance funcional)}

Repetir el mismo estado ($\theta$, distancia, m\'as adelante pose/mapa) por
Wi-Fi o Bluetooth del ESP32. Actualizar este archivo al empezar.

\subsection{Paso 3 --- Cuando Federico tenga ruedas/AS5600/DRV8833}

Contrato de capa~2 de encoders y motores; pose $(x,y,\theta)$; mapa V1
(grilla). No adelantar stubs vac\'ios en el repo de Santiago.

\subsection{Fuera de este plan de grupo}

Intentos de pose solo con IMU (integrar aceler\'ometro, trayectoria a mano
como laboratorio personal) \textbf{no} forman parte de este informe.

%==============================================================================
\section{Registro de avances}
%==============================================================================

\begin{longtable}{>{\raggedright\arraybackslash}p{0.18\textwidth}
                  >{\raggedright\arraybackslash}p{0.22\textwidth}
                  >{\raggedright\arraybackslash}p{0.52\textwidth}}
\textbf{Fecha} & \textbf{Qu\'e} & \textbf{Notas} \\
\hline
2026-10-03 & Creaci\'on de este informe & Congela decisiones de plataforma,
scan, capas, AS5600/DRV8833, V0/V1, protocolo, 3 tareas. MAIN congelado.
CONTEXTO absorbido. C\'odigo ESP32 pendiente de alineaci\'on al momento de crear
este registro. \\
2026-10-03 & Alineaci\'on firmware ESP32 V0 & HAL partida, init independiente,
tres tareas, unidades en metros, tests aislados y Processing validado.
Build PlatformIO exitoso en \texttt{esp32dev}, \texttt{imu\_only} y
\texttt{tof\_only}; carga con exit code 0. \\
2026-10-03 & Organizaci\'on del workspace & Documentos en \ruta{docs/};
\ruta{shared/} con config y tres capas; filtros PlatformIO/CIAA alineados. \\
2026-10-03 & Port inicial EDU-CIAA & HAL separada de tiempo, IMU y ToF;
drivers compartidos y CSV V0. Etapa previa a activar FreeRTOS en CIAA. \\
2026-10-03 & Runtime FreeRTOS com\'un & Tareas adquirir / drivers / publicar
movidas a \ruta{shared/layer3\_app}; ambos adaptadores crean el runtime.
Build PlatformIO y build limpio EDU-CIAA enlazan; falta prueba en placa. \\
\end{longtable}

Plantilla para filas nuevas: fecha, m\'odulo tocado, si cambi\'o un contrato
(s\'i/no), c\'omo probarlo.

%==============================================================================
\section{Checklist al cerrar un cambio}
%==============================================================================

\begin{itemize}
  \item ¿Sigue valiendo la tabla de decisiones? Si no, editar este .tex
        \textbf{en el mismo} cambio.
  \item ¿MAIN queda contradicho? Agregar o ajustar un \emph{desv\'io}, no
        editar MAIN.
  \item ¿Processing / serie rotos? Actualizar el parser o el firmware,
        un solo contrato V0.
  \item ¿Se puede probar IMU sin ToF? Si se toc\'o HAL, la respuesta debe
        seguir siendo s\'i.
\end{itemize}

\end{document}
